\section{The actual normalized theta uniformization and adelic dictionary}
\label{sec:actual-normalized-member}

The remaining complex supplier for the good arithmetic cover is an
identification of the \emph{original} primitive theta line with the
analytic theta function, including its scalar.  A polarization alone
would not determine this identification.  Here the primitive graph
identity fixes its semicharacter.  The integral coefficient action
then removes the squared theta multiplier, and a Gaussian integral
fixes the probability normalization exactly.

Retain the original host, datum, rational Hermitian space $V$, and
coefficient embeddings.  Write $F=E^+$, $r=p+q\ge1$, and
$K_0=\mathbb Q(\sqrt{-d})$, where $d\equiv7\pmod {16}$ is the
original prime.  Put $w=(1+\sqrt{-d})/2$, so $w+\bar w=1$.
The extension $E/F$ is unramified at every finite place.  In this
section $A=1$, exactly as required by the deep-cover consumer, and
$M\ge3$ is its original odd half-level.  In particular the half-ideal
$\mathfrak b$ is extended from $K_0$ and
$\mathfrak b\bar{\mathfrak b}=M\mathcal O_E$,
$\mathfrak b+\bar{\mathfrak b}=\mathcal O_E$.

The original lattice at the identity is
\begin{equation}\label{eq:member-original-lattice}
 \Lambda_*=\mathfrak a L_0,\qquad
 \mathfrak a=\texttt{AA},\qquad
 (\xi_{\rm host})\mathfrak a\bar{\mathfrak a}
                       =\mathfrak D_{E/\mathbb Q}^{-1}.
\end{equation}
Thus $\Lambda_*$, not $L_0$, is trace self-dual.  We keep the exact
original embedding $\iota:\mathrm{La}\simeq\Lambda_*$, its trace
pairing, and the source coordinate conventions.  In the actual
deep-cover datum, $\iota$ is the restriction of the given rational
frame-coordinate map to $\mathfrak a L_0$; its image and trace
isometry follow directly from the host identity above, as in
\texttt{ExistsCtxOkOfPins}, lines 177--206.  Thus the embedding is
constructed from the prescribed data, not an additional assumption.
The coordinate conventions are
\[
                 G=-V.\mathrm{gram}^{\,t},
                 \qquad \xi=-\xi_{\rm host}.
\]
The form $h_G(v,v')=v^tG\bar v'$ is linear in its first argument.
For $\tau$ in the prescribed CM type $\Psi$, it has $q$ positive
directions followed by $p$ negative directions, and
$\operatorname{Im}\tau(\xi)<0$.  The alternating pairing on every
component lattice is
\[
                 B(v,v')=\operatorname{Tr}_{E/\mathbb Q}
                                      (-\xi h_G(v,v')).
\]

\begin{theorem}\label{thm:actual-normalized-member}
For every good complex embedding of the original coefficient field,
the original $A=1$ PEL moduli scheme admits the following suppliers.
\begin{enumerate}
\item There is an actual finite \texttt{FamilyUnif}, actual
\texttt{Data} models in the fixed rational space $V$, and an
\texttt{IsUnf} map to its original framed points.  The characters
$\chi$ are identically one.  The original period formulas, tangent
frames, half-level labels, primitive line, rigidification and cube
structure are preserved.
\item The weight-reading assertion holds for every original
$\texttt{Dom}(a,n)$: $a$ is antitone separately at each embedding
and $n$ is even.  At every frame whose leading Gram minors all have
absolute value one, the theta section has exactly probability
$L^2$ norm one in the source metric.
\item There is an actual adelic dictionary satisfying all six
\texttt{DictIsUnif} clauses, the original \texttt{DictLat} marking,
\texttt{ModelW}, and the common orthogonal-base-frame clause.
Its compact subgroup is precisely the prescribed principal level
$K(M)$ on $L_0$; its lattice is \eqref{eq:member-original-lattice}.
Every dictionary frame has all Gram minors of absolute value one.
\end{enumerate}
No field extension or scalar relaxation is introduced by these complex
constructions.  Together with the actual level cover and component
field of Sections~\ref{sec:actual-level-cover} and
\ref{sec:actual-component-field}, these suppliers give the original
\texttt{Deep.Good} cover with its stated field-size budget.
\end{theorem}

\begin{proof}
\par\medskip\noindent\textbf{The actual component lattices and families.}
The Riemann correspondence identifies a complex object with its
integral homology lattice, its coefficient action, the exact
alternating pairing, its half-level labels and its positive complex
structure\cite{member-milne-av}.  The actual primitive line is
$L=(1,w)^*\mathcal P_\lambda$; hence these data determine also the
line and its rigidified cube structure.  The original theta
commutator is
$\exp(-2\pi i B(v,v')/N)$ on the classes $v/N,v'/N$,
with the source sign, by the Poincare pairing
calculation\cite{member-lan}.

The local self-dual-lattice calculation in
Section~\ref{sec:actual-component-field} applies with the fixed $\xi$:
all these rational Hermitian spaces are \emph{isometric} to $V$,
and their lattices are locally isometric to $\Lambda_*$.  At a
prime dividing $M$, $K_0$ is split.  Indeed an inert or ramified
prime of $K_0$ would be fixed by conjugation and could not occur
in two coprime complementary ideals extended from $K_0$.
At such a place the unitary lattice group is a general linear
group over the local integer ring.  Its reduction onto every
finite congruence quotient is surjective.  Complete the half-label
using the fixed root $\exp(-2\pi i/M)$ and the perfect pairing,
as in the half/full-level construction.  Consequently the local
isometries can be chosen to preserve exactly these labels.
There is no approximation assertion for the unitary group here.

Choose representatives $(\Lambda_i,t_i)$ of the resulting labelled
integral isometry classes.  Their period domains are connected:
the positive/negative decomposition domain is a product of
unitary symmetric domains, and choices of complex eigenbases
form connected general linear groups.  Local analytic triviality
of homology and the local period description make each labelled
class an open and closed analytic component of the actual moduli
scheme.  There are finitely many components, since this scheme
is of finite type over $\mathbb C$.  It is nonempty by the original
datum's positive complex structure.  This proves the required
finite nonempty indexing, including definite signatures.

Here $t_i:\mathrm{La}\to\Lambda_i$ is an
$\mathcal O_E$-linear lift of the full label modulo $M$.
Such a lift exists because $\mathrm{La}$ is projective.  It is an
isomorphism modulo $M$, precisely the two \texttt{t\_inj},
\texttt{t\_surj} clauses.  For every positive period map $P$,
take $\mathbb C^g/P\Lambda_i$, with its original labels, coefficient
action and primitive line.  Riemann's theorem algebraizes this
polarized torus and its homomorphisms.  The fine PEL moduli property
gives its classifying point and exponential.  Their kernel, coefficient
action, all-$N$ Weil pairings and level labels are exactly the
corresponding \texttt{FamilyUnif} clauses.

This construction is holomorphic in every holomorphic family of
periods: the quotient is a proper holomorphic family of polarized
tori with fixed discrete labels, hence its classifying map to the
analytic fine moduli space is holomorphic.  Conversely the integral
homology local system of a pulled-back family is constant on a
small ball, producing the stipulated local period lift.  Differentiating
the exponential gives the stipulated unique tangent vector in
each direction.  The isomorphism criterion for polarized tori
gives exactly the source fibre criterion, with the half-label
and $A=1$ congruences.  These are statements about the actual
moduli family, rather than a new quotient carrying similar points.

\par\medskip\noindent\textbf{The arithmetic frame data.}
For each representative put
\[
 \Gamma_i=\{\gamma\in U(G)(\mathbb Q):
       \gamma\Lambda_i=\Lambda_i,\quad
       (\gamma-1)\Lambda_i\subset M\Lambda_i\}.
\]
This is exactly the original half-label stabilizer: preservation
of the pairing and the adjoint coefficient action recover the
missing half of the label.  It has finite index in the lattice
stabilizer by reduction modulo $M$.  A common nonzero integer
clears the values of $h_G$ on $\Lambda_i$, giving the original
\texttt{den} field.  Choose $\kappa_0=2\xi$ and an integer
$m_i>0$ such that
$m_i\kappa_0h_G(\Lambda_i,\Lambda_i)\subset M\mathcal O_E$.
For an isotropic $v\in\Lambda_i$ and an integral $s\in F$,
\[
     u\longmapsto u+m_i s\kappa_0 h_G(u,v)v
\]
is unitary, has the same expression with a minus sign as inverse,
and is congruent to the identity modulo $M$.  This supplies the
full original transvection clause with one $m_i$, for all $v,s$.

If $\gamma\in\Gamma_i$ preserves a definite complete flag, its
matrix in an orthogonal flag basis is diagonal unitary.  This is
true at every infinite place.  Each eigenvalue $\lambda$ of its
integral $\mathbb Z$-matrix consequently has all conjugates of
absolute value one, while $(\lambda-1)/M$ is an algebraic integer.
If $\lambda\ne1$, its norm has absolute value at least one and
strictly less than one, since $2/M<1$.  Thus every eigenvalue is
one; diagonalizability gives $\gamma=1$.  This is precisely
\texttt{Gamma\_free}.  Take $\chi=1$ for every weight.  All remaining
\texttt{Data} fields follow from the given signature, lattice and
coefficient action.

\par\medskip\noindent\textbf{The primitive line selects the characteristic.}
In Appell--Humbert notation the semicharacter of the Poincare
biextension on $A\times A$ is
\[
       \alpha_{\mathcal P}(v,v')=(-1)^{B(v,v')}.
\]
Indeed it is trivial on both axes by their rigidifications, and
the semicharacter identity on the sum of the two axes gives this
formula.  Pulling back by $(1,w)$ gives, with no translation,
\begin{equation}\label{eq:member-semicharacter}
       \alpha_L(v)=(-1)^{B(v,wv)}.
\end{equation}
Its alternating form is $B$, because the adjoint of $w$ is
$\bar w$ and $w+\bar w=1$.  The Appell--Humbert classification
therefore identifies this very rigidified line, not just its
polarization\cite{member-milne-cm}.

Write $d=16k+7$ and $e=w^4\bmod8$.  The equation
$w^2-w=-(4k+2)$ gives
\[
       e^2=e,\qquad e+\bar e=1,\qquad e\bmod2=w\bmod2.
\]
The two summands $e(\Lambda_i/8\Lambda_i)$ and
$(1-e)(\Lambda_i/8\Lambda_i)$ are complementary Lagrangians:
the adjoint of $e$ is $1-e$, so each is isotropic, and perfectness
pairs them perfectly.  They are free of rank $g$ over
$\mathbb Z/8$.  Choose a symplectic basis adapted to them and
lift it to an integral symplectic basis.  Such a lift exists
because elementary symplectic row operations generate
$\mathrm{Sp}_{2g}(\mathbb Z/8)$; lifting their parameters to
$\mathbb Z$ proves surjectivity of reduction.  Choose the
source sign $B=-\omega_{\mathrm{std}}$ for this basis, denoted
$(E_1,\ldots,E_g,F_1,\ldots,F_g)$.
Formula~\eqref{eq:member-semicharacter} is then
\[
              \alpha_L(mE+nF)=(-1)^{m^tn}.
\]
Thus the original primitive line has \emph{zero} theta
characteristic in this adapted basis.  For any other fixed
symplectic basis its half-characteristic is still constant in
the period parameter; this also supplies the unrestricted
basis clause in \texttt{FamilyUnif.theta}.

\par\medskip\noindent\textbf{The exact theta normalization.}
Let $x=(x_\tau)_{\tau\in\Psi}$ be in the original domain
$GG(\Psi,G,q)$, whose leading Gram minors have the signs of
$\operatorname{diag}(1_q,-1_p)$.  Use exactly the source periods:
the $p$ coordinates of type $\tau\in\Psi$ are the last $p$
coordinates of $x_\tau^{-1}\tau(v)$; the $q$ coordinates of type
$\bar\tau$ are
$h_{G,\tau}(x_\tau e_{q-1-i},\tau(v))$, in the specified reversed
order.  They depend holomorphically on $x$ because $h_G$ is
linear in its first argument.  Indefinite Gram--Schmidt proves
that these periods are injective, discrete and span over
$\mathbb R$, with the positive Riemann form $B$.

Let $h=h(x)$ and $\tau=\tau(x)$ normalize the adapted periods:
\[
       hP(E_j)=e_j,\qquad hP(F_j)=\tau e_j,\qquad
       \tau=\tau^t,\quad Y=\operatorname{Im}\tau>0.
\]
The first period matrix is invertible by the Riemann relations,
so $h,\tau$ are holomorphic on the whole frame domain.
For the original ordered tangent types $\mathrm{ty}_1,\ldots,
\mathrm{ty}_g$ put
\begin{equation}\label{eq:member-kappa}
 P_\xi=\prod_{\nu=1}^g2|\mathrm{ty}_\nu(\xi)|,\qquad
 C_\xi=\sqrt{2^g/P_\xi},\qquad
                       \kappa(x)^2=C_\xi\det h(x).
\end{equation}
Choose either square root at each point.  The theta-frame
section is $\kappa(x)\theta(h(x)z,\tau(x))$, with
\[
       \theta(u,\tau)=\sum_{n\in\mathbb Z^g}
                         \exp(\pi i n^t\tau n+2\pi i n^tu).
\]
Its usual factor of automorphy and
\eqref{eq:member-semicharacter} identify it with a section of
the actual primitive line, through the comparison equal to one
at the rigidification.  It is nonzero: its Fourier coefficient
of index zero is one.  The same Fourier recurrence shows that
the section space of this principal line is one dimensional.
For joint holomorphy use the \emph{standard theta factors}
in $(x,z)$, which are holomorphic in $h,\tau,z$.
They define a holomorphic rigidified line on the proper
family of tori.  Its difference from the original primitive
line is trivial on every fibre by the preceding semicharacter
calculation.  Relative analytic seesaw identifies this difference
with a line from the base: proper coherent base change
\cite[\S7, Satz~5]{member-grauert}, with
constant fibre dimension $h^0=1$, makes its direct image a
line and the evaluation map an isomorphism.  The zero
rigidification trivializes that line.  Thus the normalized
comparison is jointly holomorphic and its local instances agree.

Together with the exponential derivatives this defines the
required original framed point.  In the Appell--Humbert frame
its theta function is
\[
 \Theta_x(z)=\kappa(x)e^{q_x(z)}\theta(hz,\tau),\qquad
 q_x(z)=\frac{\pi}{2}(hz)^tY^{-1}(hz),\qquad
 H_x=h^tY^{-1}\bar h.
\]
The line frame is multiplied by $e^{-q_x}$, which equals one
at zero.  Its transition functions, rigidification and
semicharacter automorphy are consequently exactly the
three original frame clauses.  These metric frames are used
only at a fixed $x$; $Y^{-1}$ need not depend holomorphically
on $x$.  Joint holomorphy uses the standard theta frame just
constructed, rather than this fibrewise metric gauge.

If all leading Gram minors have absolute value one, Gram--Schmidt
writes $x=x_0b$ with $x_0$ orthogonal and $b$ upper unipotent.
At $x_0$ the above period coordinates give the diagonal
Riemann Hermitian matrix with entries
$2|\mathrm{ty}_\nu(\xi)|$.  The coordinate change induced by
$b$ has determinant one.  Hence
\begin{equation}\label{eq:member-hermitian-determinant}
           \det H_x=P_\xi,\qquad
           \det H_x\det Y=|\det h|^2,\qquad
           |\kappa|^2=\sqrt{2^g\det Y}.
\end{equation}
On the half-open fundamental parallelepiped
$z=h^{-1}(u+\tau v)$, $u,v\in[0,1)^g$, the metric factor is
\[
 |\Theta_x(z)|^2e^{-\pi H_x(z,z)}
     =|\kappa|^2|\theta(u+\tau v,\tau)|^2
                         e^{-2\pi v^tYv}.
\]
Integrating first in $u$ kills all unequal Fourier indices.
The remaining sum over $n$ and integral over $v$ tile
$\mathbb R^g$, and give
\begin{equation}\label{eq:member-mean}
 \operatorname{mean}_{\mathbb C^g/P\Lambda_i}
        |\Theta_x(z)|^2e^{-\pi H_x(z,z)}
       =|\kappa|^2\int_{\mathbb R^g}e^{-2\pi t^tYt}\,dt
       =\frac{|\kappa|^2}{\sqrt{\det(2Y)}}=1.
\end{equation}
Absolute convergence, uniformly on the compact parallelepiped,
justifies the Fourier integration; Tonelli applies to the
remaining nonnegative sum.  The real Jacobian is the volume
of that parallelepiped and cancels in the \emph{mean}.
Thus \eqref{eq:member-mean} is the original quotient by its
Lebesgue volume, not a changed Haar normalization.

\par\medskip\noindent\textbf{Removal of the arithmetic multiplier.}
For $\gamma\in\Gamma_i$, let $\delta=\det_E\gamma$.
It is an algebraic unit, $\delta-1\in M\mathcal O_E$, and
$|\sigma(\delta)|=1$ for every complex embedding.  If $\delta\ne1$,
\[
 M^{[E:\mathbb Q]}
     \le |N_{E/\mathbb Q}(\delta-1)|
     \le 2^{[E:\mathbb Q]},
\]
a contradiction.  Thus $\det_E\gamma=1$.
Its integral symplectic matrix commutes with $e$, so in the
adapted basis its off-diagonal blocks are divisible by $8$.
Moreover
\[
                 \det_{\mathbb Z}(e\gamma+1-e)\equiv1\pmod8.
\]
Here $e$ denotes the integral endomorphism $w^4$ before reduction:
modulo $8$, idempotence gives
$\det_E(e\gamma+1-e)=e\det_E\gamma+1-e=1$;
restriction of scalars and the field norm give the displayed
integral determinant.  The same endomorphism is
$\operatorname{diag}(A,I)$ modulo $8$, so $\det A\equiv1\pmod8$.
These statements apply equally to $\gamma^{-1}$ and to the
transpose block convention for changing periods.

An element of $\mathrm{SL}_g(\mathbb Z)$ lifts $A\bmod8$:
for $g\ge2$ use elementary generation over the local ring
$\mathbb Z/8$, and for $g=1$ use $1$.  Factoring off its
block-diagonal symplectic lift leaves an element of
$\Gamma_g(8)$.  The squared zero-characteristic theta
transformation on $\Gamma_g(8)\subset\Gamma_g(4,8)$ is
\begin{equation}\label{eq:member-theta-square}
 \theta(j^{-t}u,(A\tau+B)(C\tau+D)^{-1})^2
       =\det(j)\,e^{2\pi i u^tj^{-1}Cu}\theta(u,\tau)^2,
       \qquad j=C\tau+D.
\end{equation}
The classical theta-constant transformation gives its squared
scalar exactly one\cite{member-theta-transformation}.  To pass
from constants to all $u$, first use the \emph{unsquared} function
\[
 e^{-\pi i u^tj^{-1}Cu}
       \theta(j^{-t}u,(A\tau+B)(C\tau+D)^{-1}).
\]
Substitution of the two lattice translations shows that this
function has the same factors as $\theta(u,\tau)$.  The
one-dimensional Fourier recurrence therefore makes their ratio
constant in $u$.  On the open set where $\theta(0,\tau)\ne0$ the
theta-constant formula determines the \emph{square} of this
constant as $\det j$.  The recurrence is applied to $L$,
not to the larger section space of $L^2$.
Analytic continuation gives
the identity for every $\tau$.  The block
$\operatorname{diag}(U,U^{-t})$, $\det U=1$, has exact theta
transformation by reindexing $n\mapsto U^tn$.  Composition
therefore gives \eqref{eq:member-theta-square} for the actual
arithmetic period change, with no omitted root of unity.

For $x'=\gamma x$ the period normalization is
$h'=j^{-t}h$, hence
$\kappa'^2=\kappa^2/\det j$.  The comparison of the two
actual primitive lines is rigidified.  Combining its theta
factor with \eqref{eq:member-theta-square} gives for the
section ratio $t_0$
\[
                          t_0^2
                          =\kappa'^2\det j/\kappa^2=1.
\]
This argument does not evaluate a nonzero theta section at
zero.  Equivalently one cancels the entire function $\theta^2$
on its dense nonzero locus and then evaluates the resulting
entire comparison at zero, where its rigidified value is one.
It is valid also on the theta-null divisor.

\par\medskip\noindent\textbf{The exact Borel weights and holomorphy.}
For $x'=xb$, with $b$ upper triangular, write its diagonal
blocks of sizes $q,p$ as $b_{11},b_{22}$.  The exact period
coordinate change, in the source type order, is
\[
        h_0=\operatorname{diag}(b_{11}^{\,t},b_{22}^{-1}),
 \qquad \det h_0=\frac{\det b_{11}}{\det b_{22}}.
\]
The first block is reordered in the same reversed order as
the conjugate periods.  It does not change this determinant.
We have $P'=h_0P$, $h'=hh_0^{-1}$, $\tau'=\tau$, and the
rigidified section comparison gives
$t_0^2=(\det h_0)^{-1}$.
The diagonal tangent changes are $b_{jj}^{-1}$ for $j<q$
and $b_{jj}$ for $j\ge q$.  Therefore every original even
$n$ gives exactly the source exponents
\[
 \lambda_{\tau j}=
 \begin{cases}
  -a(\bar\tau,q-1-j)-n/2,&j<q,\\
   a(\tau,j-q)+n/2,&j\ge q ,
 \end{cases}
\]
with $\tau,\bar\tau$ relabelled through the given coefficient
embedding as in \texttt{laU}.  For $\Gamma_i$ the tangent
change is the identity and $t_0^n=1$.  These prove the two
transformation clauses with the original character $\chi=1$.

Locally $\det h$ has a holomorphic square root.  Two such
roots differ by a sign, which every even $n$ removes.
Thus the chosen pointwise $\kappa$ gives holomorphic readings
for \emph{all} original $\texttt{Dom}(a,n)$.
Algebraic readings are holomorphic here: on an affine chart
of the actual framed moduli torsor, a natural reading on the
original finite type reduced test algebras is a regular function
by Yoneda.  The tangent and principal-section frame torsors are
of finite type, and the already constructed base-change
property identifies these charts with the source tests.
No continuity of the arbitrary pointwise choice of sign is used.

Every original complex tangent flag is obtained by these
period coordinates: its two definite eigenflags give the
first $q$ columns and the last $p$ quotient columns of $x$;
the unused upper off-diagonal block is precisely the Borel
freedom.  The one-dimensional section space supplies the
remaining nonzero scalar.  The same linear algebra and the
polarized-torus isomorphism criterion identify two such
framed flags exactly by $\Gamma_i$ and the original Borel.
This proves the coverage and fibre clauses of \texttt{IsUnf}.
Together with \eqref{eq:member-mean} and the family above,
all its clauses have now been supplied.

\par\medskip\noindent\textbf{The original adelic marking and six dictionary clauses.}
For $g\in G(\mathbb A)$ define
\[
 \Lambda_g=V(\mathbb Q)\cap g_f\widehat{\Lambda_*}.
\]
Additive approximation, applied after putting the compact open
lattice between $N\widehat{\Lambda_*}$ and
$N^{-1}\widehat{\Lambda_*}$, proves
$\widehat{\Lambda_g}=g_f\widehat{\Lambda_*}$.
This is precisely the source \texttt{latG}, whose definition
first multiplies the unscaled adelic lattice by $\mathfrak a$;
invertible fractional ideals commute with localization,
completion and intersection here.  Thus its \texttt{latGHat}
is precisely the displayed completion.  All lattices, including
their fractional-ideal defects, are retained.

Choose a purely finite representative $g_i$ for each labelled
component.  The local isometries with their exact labels
constructed above supply it.  Choose $t_i$ so that
\[
             t_i(a)\equiv(g_i)_f\iota(a)
                       \pmod{M\widehat{\Lambda_i}} .
\]
For every $g$ choose $i$ and $\gamma\in G(\mathbb Q)$ with
$\gamma g_f\in(g_i)_fK(M)$, and set
\[
                 \operatorname{dict}(g)
                      =(i,(\gamma g)_\infty x_0).
\]
Here $x_0$ is a single orthogonal frame for the fixed $G$ at
all prescribed infinite embeddings.  Choose $\gamma=1$ at
$g=g_i$.  The original marking is now literal:
\[
 \Lambda_i=\Lambda_{\gamma g},\qquad
 \widehat{t_i(a)}-(\gamma g)_f\widehat{\iota(a)}
       \in M\widehat{\Lambda_i}
       \quad(a\in\bar{\mathfrak b}\,\mathrm{La}).
\]
This is \texttt{DictLat}, with its original rational multiplier,
selected half-submodule and source and target completions.
Principal congruence on $\Lambda_*$ is equivalent to principal
congruence on $L_0$: locally cancel the invertible ideal
$\mathfrak a$.  Hence the subgroup is exactly the source
\texttt{IsLevelK}, not a conjugated replacement.

For completeness, its six uniformization clauses are as follows.
\begin{enumerate}
\item Left multiplication by $G(\mathbb Q)$ changes only the
chosen marking; two markings of a component differ by
$\Gamma_i$.
\item Right multiplication by $K(M)$ preserves the lattice and
its selected label, and has trivial infinite component.
It therefore also changes the dictionary by $\Gamma_i$ only.
\item Each $(i,x_0)$ occurs at $g_i$.  Right multiplication by
$h_\infty$ acts on its frame through
$x_0^{-1}h_\infty x_0$, exactly the original \texttt{archMat},
up to the same change of marking.
\item Every dictionary frame is orthogonal, hence belongs to
$GG(\Psi,G,q)$ and has all leading Gram minors of absolute
value one.
\item Every frame in $GG$ is an orthogonal frame times an
upper triangular matrix.  The infinite unitary group is
transitive on orthogonal frames, so the dictionary is onto
each domain modulo $\Gamma_i$ and that Borel.
\item If two dictionary frames differ by $\Gamma_i$, the
rational marking comparison makes their infinite components
equal.  Their finite quotient preserves $\Lambda_*$ and its
half-label.  The half/full adjoint argument makes it
congruent to one modulo $M$; cancellation of $\mathfrak a$
puts it in the original $K(M)$.  Thus
$g'=\gamma gk$ for $\gamma\in G(\mathbb Q)$, $k\in K(M)$,
as required by the final clause.
\end{enumerate}
The derivative, period and lattice equations already proved
give the remaining \texttt{DictLat} clauses with the same
$x_0$, the original infinite embeddings, and exact
\texttt{ModelW}.  The actual test-lifting theorem gives
\texttt{unTest} for these framed points.  This completes the
claimed Member supplier.
\end{proof}

\begin{corollary}\label{cor:actual-good-sized-cover}
Under the original \texttt{Core.Pins}, $j\in\texttt{InR}$,
$\mathrm{splitBound}\ge5$ and $(15,g_0g_1)=1$, for every $e\ge0$
there is an actual cover $Y$ satisfying the original
\texttt{Deep.Good} contract at depth $e$ and budget
\[
  \operatorname{fieldSz}(C)+
       20\bigl(\log(g_0g_1\,15^e)+\log d+1\bigr).
\]
It has precisely $p_d=15^e$, original higher half-level
$M_2=g_0g_1\,15^e$, and $A_2=1$.
\end{corollary}

\begin{proof}
Use the actual complementary half-ideal at $15^e$, the actual
level cover, and the compositum with the $K_0$ ray field of
modulus $M_2$.  The component-field theorem gives
\texttt{FnAlgClosed}, normality over $C$, the prescribed roots,
and exactly the displayed budget.  The preceding theorem
supplies \texttt{IsUnf}, \texttt{unTest}, trivial characters,
the six dictionary laws, unit flag minors, \texttt{ModelW},
and the original common orthogonal marking.  The retained
host identity for $|\tau(\xi)|$ is the original
\texttt{xiAbs} clause.  These are precisely all fields of
\texttt{Deep.Good}; neither an automorphic-character supplier
nor an enlarged field-size bound is being inserted.
\end{proof}

The corollary closes the actual good-cover supplier, rather than a
new abstract sufficient condition.  The global arithmetic section
estimate, signed boundary comparison and the resulting W10 degree
bound remain separate obligations.
